\section*{Chapitre \ref{ch:b1:weathering-clays} --- L'altération et les argiles}

\begin{solution}{exo:b1:weathering-clays:1}
\ce{K+}~: $1/1.38 = 0.72$~; \ce{Fe^2+}~: $2/0.78 = 2.56$~; tous deux inférieurs à 3, cations
solubles. \ce{Fe^3+}~: $3/0.645 = 4.65$~; \ce{Ti^4+}~: $4/0.605 = 6.61$~; tous deux
entre 3 et 10, hydrolysats. Dans un sol gorgé d'eau, sans oxygène, le fer est
libéré sous forme de \ce{Fe^2+}, qui reste en solution et est emporté par
l'eau~; dans un sol bien drainé, l'oxygène le transforme en \ce{Fe^3+}, qui
précipite aussitôt en oxydes et hydroxydes et reste.
\end{solution}

\begin{solution}{exo:b1:weathering-clays:2}
A gonfle avec le glycol et s'effondre au chauffage~: une \omterm{def:b1:weathering-clays:clay-mineral}{smectite}. B est à
\SI{7}{\angstrom} et est détruite à \SI{550}{\celsius}~: la \omterm{def:b1:weathering-clays:clay-mineral}{kaolinite}. C, à
\SI{10}{\angstrom} quel que soit le traitement~: l'\omterm{def:b1:weathering-clays:clay-mineral}{illite}. D, à \SI{14}{\angstrom}
sans gonflement ni effondrement~: la \omterm{def:b1:weathering-clays:clay-mineral}{chlorite}.
\end{solution}

\begin{solution}{exo:b1:weathering-clays:3}
Deux orthoses apportent les 2 aluminium d'une \omterm{def:b1:weathering-clays:clay-mineral}{kaolinite}~; $6 - 2 = 4$ silicium
partent sous forme de \ce{H4SiO4}~; 2 \ce{K+} partent, équilibrés par 2 \ce{H+}~; l'oxygène,
$16 + w = 9 + 16$, donne $w = 9$~:
\[
  \ce{2 KAlSi3O8 + 2 H+ + 9 H2O -> Al2Si2O5(OH)4 + 2 K+ + 4 H4SiO4}.
\]
L'hydrogène se vérifie~: $2 + 18 = 20 = 4 + 16$. Tout le potassium part, et $4/6 =
2/3$ du silicium.
\end{solution}

\begin{solution}{exo:b1:weathering-clays:4}
(a) La masse ne change pas~: $10 \times 999.87/916.72 = \qty{10.91}{cm^3}$.
(b) Entre environ \qty{-4}{\celsius} et \qty{-15}{\celsius}. Juste sous zéro
la glace ne peut accumuler assez de pression pour coincer la fissure~; dans un gel plus dur
l'eau ne se déplace plus à travers les pores vers la glace, qui cesse de
croître. (c) D'après la
\cref{prop:b1:weathering-clays:surface} avec $n = \qty{1}{m}/\qty{1}{cm} = 100$,
l'aire est multipliée par 100~: de \qty{6}{m^2} à \qty{600}{m^2}.
\end{solution}

\begin{solution}{exo:b1:weathering-clays:5}
$S/m = 6/(\rho d)$~: pour \qty{1}{cm}, $6/(2600\times0.01) =
\qty{0.231}{m^2.kg^{-1}}$~; pour \qty{1}{mm}, \qty{2.31}{m^2.kg^{-1}}~; pour
\qty{1}{\micro m}, \qty{2308}{m^2.kg^{-1}}, soit \qty{2.3}{m^2.g^{-1}}.
Les réactions chimiques ont lieu à la surface des grains, donc leur vitesse par
kilogramme de roche croît avec la surface spécifique, dix mille fois du
centimètre au micromètre. L'\omterm{def:b1:weathering-clays:weathering}{altération physique} découpe la roche en morceaux toujours
plus petits et ouvre des fissures par où l'eau entre~: elle multiplie la
surface sur laquelle l'\omterm{def:b1:weathering-clays:weathering}{altération chimique} peut agir.
\end{solution}

\begin{solution}{exo:b1:weathering-clays:6}
Une tonne d'albite fait $\num{e6}/262.22 = \qty{3813.6}{mol}$. D'après la
\cref{prop:b1:weathering-clays:feldspars}, deux moles donnent une de \omterm{def:b1:weathering-clays:clay-mineral}{kaolinite}~:
$\qty{1906.8}{mol} \times \qty{258.16}{g.mol^{-1}} = \qty{492.3}{kg}$. Deux des
trois atomes de silicium partent~: $2\times3813.6 = \qty{7627.2}{mol}$ de
\ce{SiO2}, \qty{458.2}{kg}. Tout le sodium part~: $3813.6\times22.99 =
\qty{87.7}{kg}$. Moins de la moitié de la masse du \omterm{def:b1:rock-forming-minerals:groups}{feldspath} reste sous forme d'argile.
\end{solution}

\begin{solution}{exo:b1:weathering-clays:7}
D'après la \cref{prop:b1:weathering-clays:immobile}, un kilogramme de granite laisse
$0.30/0.60 = \qty{0.5}{kg}$ de saprolite. Fraction perdue~: sodium $1 -
(0.35/3.5)(0.5) = 0.95$~; silicium $1 - (60/72)(0.5) = 0.583$~; aluminium $1 -
(27/14)(0.5) = 0.036$. L'aluminium est pratiquement conservé, le sodium a presque
entièrement disparu, et près de \qty{60}{\percent} du silicium aussi, bien qu'une grande part de ce qui
reste soit tenue dans le quartz~: les \omterm{def:b1:rock-forming-minerals:groups}{feldspaths} se transforment en \omterm{def:b1:weathering-clays:clay-mineral}{kaolinite},
\omterm{def:b1:weathering-clays:hydrolysis-degree}{monosiallitisation}.
\end{solution}

\begin{solution}{exo:b1:weathering-clays:8}
Pyrophyllite~: $x = y = 0$, charge 0. Feuillet de muscovite~: $x = 1$, charge $-1$.
Feuillet de beidellite~: $x = 0.5$, charge $-0.5$. Feuillet de montmorillonite~: $y = 0.33$,
charge $-0.33$. Les deux feuillets de \omterm{def:b1:weathering-clays:clay-mineral}{smectite}, beidellite et montmorillonite, peuvent
gonfler~: leur petite charge est équilibrée par peu de cations interfoliaires, qui tiennent les
feuillets lâchement et attirent l'eau en s'hydratant. La pyrophyllite n'a pas de
cations interfoliaires à hydrater, et la charge $-1$ de la muscovite est équilibrée par des
ions potassium qui verrouillent les feuillets.
\end{solution}

\begin{solution}{exo:b1:weathering-clays:9}
(a) Mg~: 2 et 2~; Si~: 1 et 1~; C~: 4 et 4~; O~: $4 + 8 + 4 = 16$ et $12 + 4 =
16$~; H~: 8 et $4 + 4 = 8$~; charge~: 0 et $4 - 4 = 0$. (b) La masse molaire de
la forstérite est $2\times24.305 + 28.085 + 4\times15.999 =
\qty{140.69}{g.mol^{-1}}$, donc une tonne fait \qty{7108}{mol}. Elle consomme
$4\times7108 = \qty{28431}{mol}$ de \ce{CO2}, \qty{1251}{kg}~; avec deux
magnésium par unité formulaire, la \cref{prop:b1:weathering-clays:co2} en stocke
$2\times7108$ moles, \qty{626}{kg}, la moitié de la quantité consommée. (c) Une tonne de
calcite, \qty{9991}{mol}, consomme \qty{9991}{mol} de \ce{CO2} (\qty{440}{kg})
en se dissolvant, et le rend en totalité quand le carbonate précipite de nouveau~:
rien n'est stocké.
\end{solution}

\begin{solution}{exo:b1:weathering-clays:10}
(a) $3 < 3/r < 10$ donne $\qty{0.3}{\angstrom} < r < \qty{1.0}{\angstrom}$.
(b) $3/0.27 = 11.1 > 10$~: le bore voyage sous forme d'oxyanion, le borate. (c) $2/0.83 =
2.41$, un cation soluble~; $4/0.53 = 7.55$, un hydrolysat. Dans une eau pauvre en
oxygène le manganèse est dissous sous forme de \ce{Mn^2+}~; là où la source rencontre l'air,
il est oxydé en \ce{Mn^4+}, qui précipite aussitôt en un oxyde noir sur
la roche. (d) Son \omterm{def:b1:weathering-clays:ionic-potential}{potentiel ionique}, $4/0.40 = 10.0$, se trouve sur la limite
entre hydrolysats et oxyanions~: le silicium n'est ni retenu comme l'aluminium ni
emporté comme le soufre, et la fraction lessivée dépend de la quantité d'eau qui
passe, la base des trois degrés d'\omterm{def:b1:weathering-clays:hydrolysis}{hydrolyse}.
\end{solution}

\begin{solution}{exo:b1:weathering-clays:11}
(a) Avec \qty{12.011}{g.mol^{-1}} de carbone, \num{11.7e12} à
\qty{13.6e12}{mol} font 0.141 à 0.163 milliard de tonnes de carbone par an~; avec
\qty{44.009}{g.mol^{-1}}, 0.515 à 0.599 milliard de tonnes de dioxyde de carbone.
(b) $4.08/13.6 = \qty{30}{\percent}$ à $4.08/11.7 = \qty{35}{\percent}$.
(c) La moitié du dioxyde de carbone consommé est stockée quand le calcium et le magnésium
précipitent sous forme de carbonate~: 0.070 à 0.082 milliard de tonnes de carbone par an. Le
chiffre réel est plus petit, parce qu'une partie du bicarbonate vient du sodium et du
potassium, qui ne forment pas de carbonate sur le fond marin.
\end{solution}

\begin{solution}{exo:b1:weathering-clays:12}
Si/Al $= (w_{\ce{SiO2}}/60.08)/(2w_{\ce{Al2O3}}/101.96)$. A~: $0.7748/0.7748 =
1.00$, le rapport de la \omterm{def:b1:weathering-clays:clay-mineral}{kaolinite}, \omterm{def:b1:weathering-clays:hydrolysis-degree}{monosiallitisation}. B~: $0.9987/0.4904 = 2.04$, un
\omterm{def:b1:weathering-clays:clay-mineral}{minéral argileux} 2:1, \omterm{def:b1:weathering-clays:hydrolysis-degree}{bisiallitisation}. C~: $0.0333/1.2358 = 0.027$, presque plus
de silicium, gibbsite, \omterm{def:b1:weathering-clays:hydrolysis-degree}{allitisation}. Du lessivage le moins intense au plus intense,
c'est-à-dire vers le climat le plus chaud et le plus humide~: B, A, C.
\end{solution}

\begin{solution}{pb:b1:weathering-clays:1}
\textbf{1.} \ce{2 NaAlSi3O8 + 2 H+ + 9 H2O -> Al2Si2O5(OH)4 + 2 Na+ + 4 H4SiO4}.

\textbf{2.} \ce{CaAl2Si2O8 + 2 H+ + H2O -> Al2Si2O5(OH)4 + Ca^2+}. Aucun silicium
ne part~: l'anorthite a déjà le silicium par aluminium de la \omterm{def:b1:weathering-clays:clay-mineral}{kaolinite}.

\textbf{3.} Une gibbsite prend un aluminium, soit une demi-anorthite~: en doublant,
\ce{CaAl2Si2O8 + 2 H+ + 6 H2O -> 2 Al(OH)3 + Ca^2+ + 2 H4SiO4}. Oxygène~:
$8 + 6 = 6 + 8$~; hydrogène~: $2 + 12 = 6 + 8$.

\textbf{4.} \ce{5 NaAlSi3O8 + 4 H+ + 16 H2O -> 2 Na_{0.5}Al2(Si_{3.5}Al_{0.5})O10(OH)2 + 4 Na+ + 8 H4SiO4}; % ce-unbalanced-ok: decimal subscripts the checker cannot read; balance tested in tests/bachelor-1/test_weathering_clays.py
hydrogène, $4 + 32 = 36$ à gauche et $2\times2 + 8\times4 = 36$ à
droite, équilibré.

\textbf{5.} Chaque \ce{H+} vient avec un \ce{HCO3-} issu d'un \ce{CO2} et d'un
\ce{H2O}~:
\ce{2 NaAlSi3O8 + 2 CO2 + 11 H2O -> Al2Si2O5(OH)4 + 2 Na+ + 2 HCO3- + 4 H4SiO4}
et \ce{CaAl2Si2O8 + 2 CO2 + 3 H2O -> Al2Si2O5(OH)4 + Ca^2+ + 2 HCO3-}.

\textbf{6.} L'albite en beidellite est de la \omterm{def:b1:weathering-clays:hydrolysis-degree}{bisiallitisation}~; l'albite et l'anorthite en
\omterm{def:b1:weathering-clays:clay-mineral}{kaolinite} sont de la \omterm{def:b1:weathering-clays:hydrolysis-degree}{monosiallitisation}~; l'anorthite en gibbsite est de l'\omterm{def:b1:weathering-clays:hydrolysis-degree}{allitisation}.

\textbf{7.} $\num{e6}/262.22 = \qty{3813.6}{mol}$ d'albite~;
$\num{e6}/278.20 = \qty{3594.5}{mol}$ d'anorthite.

\textbf{8.} Albite~: $3813.6/2 = \qty{1906.8}{mol}$ de \omterm{def:b1:weathering-clays:clay-mineral}{kaolinite},
\qty{492.3}{kg}. Anorthite~: \qty{3594.5}{mol}, \qty{928.0}{kg}.

\textbf{9.} Albite~: $2\times3813.6 = \qty{7627.2}{mol}$ de \ce{SiO2},
\qty{458.2}{kg}, et \qty{3813.6}{mol} de sodium, \qty{87.7}{kg}. Anorthite~:
pas de silice, et \qty{3594.5}{mol} de calcium, \qty{144.1}{kg}.

\textbf{10.} Albite~: $1/2.62 = \qty{0.382}{m^3}$~; sa \omterm{def:b1:weathering-clays:clay-mineral}{kaolinite},
$0.4923/2.6 = \qty{0.189}{m^3}$~: le solide perd \qty{50.4}{\percent} de son
volume. Anorthite~: $1/2.75 = \qty{0.364}{m^3}$~; sa \omterm{def:b1:weathering-clays:clay-mineral}{kaolinite},
$0.9280/2.6 = \qty{0.357}{m^3}$~: une perte de seulement \qty{1.8}{\percent}.

\textbf{11.} L'espace laissé devient des pores, remplis d'eau et d'air~: une
saprolite sur un granite riche en albite est légère, poreuse et friable, ses grains de quartz
libres dans un peu d'argile. L'anorthite se transforme en presque son propre volume de
\omterm{def:b1:weathering-clays:clay-mineral}{kaolinite}, donc le gabbro altéré garde l'essentiel de sa matière solide.

\textbf{12.} Le sodium ($z/r = 0.98$) et le calcium (2.0) sont des cations solubles et
partent~; l'aluminium (5.6) est un hydrolysat et reste entièrement dans la \omterm{def:b1:weathering-clays:clay-mineral}{kaolinite}~;
le silicium (10.0), sur la limite, part en partie de l'albite et pas du tout
de l'anorthite, dont la \omterm{def:b1:weathering-clays:clay-mineral}{kaolinite} peut retenir le silicium.

\textbf{13.} Un \ce{CO2} par albite~: \qty{3813.6}{mol}, $3813.6 \times 44.01 =
\qty{167.8}{kg}$.

\textbf{14.} Deux par anorthite~: \qty{7189.0}{mol}, \qty{316.4}{kg}.

\textbf{15.} $7189.0 \times 61.02 = \qty{438.7}{kg}$ de bicarbonate.

\textbf{16.} Des organismes le précipitent sous forme de carbonate de calcium,
\ce{Ca^2+ + 2 HCO3- -> CaCO3 + CO2 + H2O}~: pour chaque calcium, l'un des deux
\ce{CO2} retourne à l'air, \qty{3594.5}{mol} soit \qty{158.2}{kg}, et
l'autre est enfoui dans le carbonate.

\textbf{17.} Le sodium ne forme pas de carbonate sur le fond marin~: son bicarbonate
ne peut être stocké sous forme de roche carbonatée, et le stockage à long terme se compte à partir du
calcium et du magnésium seulement.

\textbf{18.} Une tonne de calcite consomme \qty{9991}{mol} de \ce{CO2} quand elle
se dissout et le rend en totalité quand elle précipite~: rien n'est stocké, et
seule l'altération des silicates de calcium et de magnésium retire le dioxyde de carbone
de l'air pour de bon.

\textbf{19.} $\qty{200}{m^3.s^{-1}} \times \qty{0.3}{mol.m^{-3}} =
\qty{60}{mol.s^{-1}}$~; sur une année de \qty{3.156e7}{s}, \qty{1.893e9}{mol}, soit
$\num{1.893e9}/\num{50000} = \qty{37870}{mol}$ par kilomètre carré.

\textbf{20.} Deux bicarbonates par anorthite~: \qty{9.467e8}{mol} par an,
$\num{9.467e8}\times\qty{278.20}{g} = \qty{263400}{t}$ par an, soit
\qty{5.27}{t} par kilomètre carré.

\textbf{21.} La roche fraîche consommée est $5.27/0.20 = \qty{26.3}{t}$, soit
$\num{26340}/2700 = \qty{9.75}{m^3}$, par kilomètre carré et par an~: une couche
épaisse de \qty{9.75}{\micro m} par an, \qty{9.8}{m} par million d'années.

\textbf{22.} Un \ce{CO2} stocké par calcium~: $\num{9.467e8} \times 44.01 =
\qty{4.17e10}{g}$, environ \qty{41700}{t} par an.

\textbf{23.} Elle ne compte que la matière emportée en solution. La \omterm{def:b1:weathering-clays:clay-mineral}{kaolinite}
et le quartz laissés par l'altération sont aussi enlevés par l'érosion, sous forme de boue et de sable dans
la rivière, et ne sont pas comptés~; pas plus que l'altération des autres
minéraux de la roche. Et si une partie du bicarbonate vient de \omterm{def:b1:rock-forming-minerals:groups}{feldspaths} sodiques ou
potassiques, qui donnent un bicarbonate par unité formulaire au lieu de
deux, il a fallu qu'une plus grande quantité de roche s'altère pour le fournir.

\textbf{24.} D'après la \cref{prop:b1:weathering-clays:co2}, une mole de \ce{CO2} est
stockée par mole d'anorthite~: $3594.5 \times \qty{44.01}{g} = \qty{158.19}{kg}$.
L'altération d'une tonne d'anorthite retire pour de bon
\textbf{158}~kg de dioxyde de carbone.
\end{solution}
